\chapter{1994年全国卷（理）}
\section{单选题}

\begin{problem}
设全集 \(I = \{0, 1, 2, 3, 4\}\)，集合 \(A = \{0, 1, 2, 3\}\)，集合 \(B = \{2, 3, 4\}\)，则 \(\overline{A} \cup \overline{B} =\)（\quad）

\choices
    {\(\{0\}\)}
    {\(\{0,1\}\)}
    {\(\{0,1,4\}\)}
    {\(\{0,1,2,3,4\}\)}
\end{problem}

\begin{answer}
C
\end{answer}

\begin{solution}
由全集 \(I\) 中的补集定义，\(\overline{A}=\{4\}\)，\(\overline{B}=\{0,1\}\). 因此 \(\overline{A}\cup\overline{B}=\{4\}\cup\{0,1\}=\{0,1,4\}\).
故选 C.
\end{solution}

\begin{problem}
如果方程 \(x^{2} + ky^{2} = 2\) 表示焦点在 \(y\) 轴上的椭圆，那么实数 \(k\) 的取值范围是（\quad）

\choices
    {\((0, +\infty)\)}
    {\((0, 2)\)}
    {\((1, +\infty)\)}
    {\((0, 1)\)}
\end{problem}

\begin{answer}
D
\end{answer}

\begin{solution}
要表示椭圆，必须有 \(k>0\). 将方程化为 \(\frac{x^2}{2}+\frac{y^2}{2/k}=1\). 焦点在 \(y\) 轴上，说明长半轴在 \(y\) 轴方向，即 \(\frac{2}{k}>2\).
所以 \(0<k<1\)，故选 D.
\end{solution}

\begin{problem}
极坐标方程 \(\rho = \cos \left(\frac{\pi}{4} -\theta\right)\) 所表示的曲线是（\quad）

\choices
    {双曲线}
    {椭圆}
    {抛物线}
    {圆}
\end{problem}

\begin{answer}
D
\end{answer}

\begin{solution}
利用 \(x=\rho\cos\theta\)，\(y=\rho\sin\theta\)，原方程给出 \(\rho=\frac{1}{\sqrt{2}}\left(\cos\theta+\sin\theta\right)\). 两边乘以 \(\rho\)，得 \(x^2+y^2=\frac{x+y}{\sqrt{2}}\). 配方得 \(\left(x-\frac{1}{2\sqrt{2}}\right)^2+\left(y-\frac{1}{2\sqrt{2}}\right)^2=\frac14\)，
这是一个圆，故选 D.
\end{solution}

\begin{problem}
设 \(\theta\) 是第二象限的角，则必有（\quad）

\choices
    {\(\tan \frac{\theta}{2} > \cot \frac{\theta}{2}\)}
    {\(\tan \frac{\theta}{2} < \cot \frac{\theta}{2}\)}
    {\(\sin \frac{\theta}{2} > \cos \frac{\theta}{2}\)}
    {\(\sin \frac{\theta}{2} < \cos \frac{\theta}{2}\)}
\end{problem}

\begin{answer}
A
\end{answer}

\begin{solution}
因为“\(\theta\) 是第二象限的角”，所以存在 \(k\in\Z\)，使
\(\frac{\pi}{2}+2k\pi<\theta<\pi+2k\pi\). 令 \(\varphi=\frac{\theta}{2}\)，则
\(\frac{\pi}{4}+k\pi<\varphi<\frac{\pi}{2}+k\pi\).
由正切函数的周期为 \(\pi\)，得 \(\tan\varphi>1\)，从而 \(\cot\varphi=1/\tan\varphi<1\)，所以
\(\tan\frac{\theta}{2}>\cot\frac{\theta}{2}\).
另一方面，\(\varphi-\frac{\pi}{4}\in(k\pi,k\pi+\frac{\pi}{4})\)，故当 \(k\) 为偶数时 \(\sin\varphi>\cos\varphi\)，当 \(k\) 为奇数时 \(\sin\varphi<\cos\varphi\). 因此 C、D 均不恒成立，故选 A.
\end{solution}

\begin{problem}
某种细菌在培养过程中，每 \(20\) 分钟分裂一次（一个分裂为两个）. 经过 \(3\) 小时，这种细菌由 \(1\) 个可繁殖成（\quad）

\choices
    {\(511\) 个}
    {\(512\) 个}
    {\(1023\) 个}
    {\(1024\) 个}
\end{problem}

\begin{answer}
B
\end{answer}

\begin{solution}
3 小时等于 \(180\) 分钟，共分裂 \(\frac{180}{20}=9\) 次. 每次分裂后数量变为原来的 \(2\) 倍，所以最后有 \(1\cdot 2^9=512\) 个，故选 B.
\end{solution}

\begin{problem}
在下列函数中，以 \(\frac{\pi}{2}\) 为周期的函数是（\quad）

\choices
    {\(y = \sin 2x + \cos 4x\)}
    {\(y = \sin 2x\cos 4x \)}
    {\(y = \sin 2x + \cos 2x \)}
    {\(y = \sin 2x\cos 2x \)}
\end{problem}

\begin{answer}
D
\end{answer}

\begin{solution}
将 \(x\) 换成 \(x+\pi/2\) 时，\(\sin 2x\) 变为 \(-\sin 2x\)，\(\cos 2x\) 变为 \(-\cos 2x\)，而 \(\cos 4x\) 不变. 因而选项 A 的函数值变为 \(-\sin 2x+\cos 4x\)，选项 B 的函数值变为 \(-\sin 2x\cos 4x\)，选项 C 的函数值变为 \(-\sin 2x-\cos 2x\)，都不恒等于原函数值. 对于选项 D，
\(\sin 2x\cos 2x=\frac12\sin 4x\)，
而 \(\sin 4x\) 的周期为 \(\pi/2\)，故选 D.
\end{solution}

\begin{problem}
已知正六棱台的上、下底面边长分别为 \(2\) 和 \(4\)，高为 \(2\)，则其体积为（\quad）

\choices
    {\(32\sqrt{3}\)}
    {\(28\sqrt{3}\)}
    {\(24\sqrt{3}\)}
    {\(20\sqrt{3}\)}
\end{problem}

\begin{answer}
B
\end{answer}

\begin{solution}
正六边形边长为 \(s\) 时，可分成 6 个边长为 \(s\) 的正三角形，因此底面积为 \(S=6\cdot\frac{\sqrt{3}}{4}s^2=\frac{3\sqrt{3}}{2}s^2\).
上、下底面积分别为 \(6\sqrt{3}\) 和 \(24\sqrt{3}\). 棱台体积公式给出
\[
V=\frac{2}{3}\left(6\sqrt{3}+24\sqrt{3}+\sqrt{6\sqrt{3}\cdot24\sqrt{3}}\right)
=\frac{2}{3}\left(6+24+12\right)\sqrt{3}=28\sqrt{3}
\]
故选 B.
\end{solution}

\begin{problem}
设 \(F_{1}\) 和 \(F_{2}\) 为双曲线 \(\frac{x^2}{4} - y^2 = 1\) 的两个焦点，点 \(P\) 在双曲线上且满足 \(\angle F_1PF_2 = 90^\circ\)，则 \(\triangle F_1PF_2\) 的面积是（\quad）

\choices
    {\(1\)}
    {\(\frac{\sqrt{5}}{2}\)}
    {\(2\)}
    {\(\sqrt{5}\)}
\end{problem}

\begin{answer}
A
\end{answer}

\begin{solution}
双曲线的标准方程中 \(a^2=4\)，\(b^2=1\)，所以 \(c^2=a^2+b^2=5\)，两焦点间距离为 \(2c=2\sqrt{5}\). 设
\(u=PF_1\)，\(v=PF_2\).
由双曲线定义，\(\lvert u-v\rvert=2a=4\). 又因为 \(\angle F_1PF_2=90^\circ\)，勾股定理给出 \(u^2+v^2=(2\sqrt{5})^2=20\). 将 \((u-v)^2=16\) 与上式相减，得 \(2uv=4\)，即 \(uv=2\). 因此 \(S_{\triangle F_1PF_2}=\frac12uv=1\).
故选 A.
\end{solution}

\begin{problem}
如果复数 \(z\) 满足 \(|z + \i| + |z - \i| = 2\)，那么 \(|z + \i + 1|\) 的最小值是（\quad）

\choices
    {\(1\)}
    {\(\sqrt{2}\)}
    {\(2\)}
    {\(\sqrt{5}\)}
\end{problem}

\begin{answer}
A
\end{answer}

\begin{solution}
设 \(z=x+y\i\). 条件 \(|z+\i|+|z-\i|=2\) 表示点 \(z\) 到点 \((0,-1)\)、\((0,1)\) 的距离和等于两定点间距离 \(2\)，所以 \(z\) 在线段 \([-\i,\i]\) 上，即 \(x=0\)，\(-1\leqslant y\leqslant1\). 于是 \(|z+\i+1|=|1+(y+1)\i|=\sqrt{1+(y+1)^2}\geqslant1\).
当 \(y=-1\)，即 \(z=-\i\) 时取等号，故最小值为 \(1\)，选 A.
\end{solution}

\begin{problem}
有甲、乙、丙三项任务，甲需 \(2\) 人承担，乙、丙各需 \(1\) 人承担. 从 \(10\) 人中选派 \(4\) 人承担这三项任务，不同的选法共有（\quad）

\choices
    {\(1260\) 种}
    {\(2025\) 种}
    {\(2520\) 种}
    {\(5040\) 种}
\end{problem}

\begin{answer}
C
\end{answer}

\begin{solution}
先从 \(10\) 人中选出承担任务的 \(4\) 人，再从这 \(4\) 人中选出 \(2\) 人承担甲，最后将剩下的 \(2\) 人分别安排给乙、丙. 因此选法数为
\[
\mathrm C_{10}^{4}\mathrm C_{4}^{2}\mathrm A_{2}^{2}=210\cdot6\cdot2=2520
\]
故选 C.
\end{solution}

\begin{problem}
对于直线 \(m\)，\(n\) 和平面 \(\alpha\)，\(\beta\)，\(\alpha \perp \beta\) 的一个充分条件是（\quad）

\choices
    {\(m \perp n\)，\(m \parallel \alpha\)，\(n \parallel \beta\)}
    {\(m \perp n\)，\(\alpha \cap \beta = m\)，\(n \subset \alpha\)}
    {\(m \parallel n\)，\(n \perp \beta\)，\(m \subset \alpha\)}
    {\(m \parallel n\)，\(m \perp \alpha\)，\(n \perp \beta\)}
\end{problem}

\begin{answer}
C
\end{answer}

\begin{solution}
由 \(n\perp\beta\)，且 \(m\parallel n\)，可知 \(m\perp\beta\). 又 \(m\subset\alpha\)，所以平面 \(\alpha\) 内有一条直线 \(m\) 垂直于平面 \(\beta\). 根据平面与平面垂直的判定定理，\(\alpha\perp\beta\)，故选 C.
\end{solution}

\begin{problem}
设函数 \(f(x) = 1 - \sqrt{1 - x^2}\)（\(-1 \leqslant x \leqslant 0\)），则函数 \(y = f^{-1}(x) \) 的图象是（\quad）

\choices
    {\choicebitmap[width=0.15\paperwidth]{img/1994/national_science/q12_optA.png}}
    {\choicebitmap[width=0.15\paperwidth]{img/1994/national_science/q12_optB.png}}
    {\choicebitmap[width=0.15\paperwidth]{img/1994/national_science/q12_optC.png}}
    {\choicebitmap[width=0.15\paperwidth]{img/1994/national_science/q12_optD.png}}
\end{problem}

\begin{answer}
B
\end{answer}

\begin{solution}
由 \(y=1-\sqrt{1-x^2}\)，\(-1\leqslant x\leqslant0\) 可得函数值域为 \([0,1]\)，且函数在该区间上严格递减. 设反函数图象上的点为 \((x,y)\)，则原函数图象上的对应点为 \((y,x)\)，从而 \(x=1-\sqrt{1-y^2}\)，\(y\leqslant0\). 因此 \(y=-\sqrt{2x-x^2}\)，\(0\leqslant x\leqslant1\).
反函数图象是第四象限内从 \((0,0)\) 到 \((1,-1)\) 的圆弧，与第二个选项的图形相符，故选 B.
\end{solution}

\begin{problem}
已知过球面上 \(A\)，\(B\)，\(C\) 三点的截面和球心的距离等于球半径的一半，且 \(AB = BC = CA = 2 \)，则球面面积是（\quad）

\choices
    {\(\frac{16}{9}\pi\)}
    {\(\frac{8}{3}\pi\)}
    {\(4\pi\)}
    {\(\frac{64}{9}\pi\)}
\end{problem}

\begin{answer}
D
\end{answer}

\begin{solution}
截面圆的圆心是球心在截面上的射影. 由于 \(ABC\) 为边长 \(2\) 的正三角形，截面圆半径为正三角形外接圆半径 \(r=\frac{2}{\sqrt{3}}\). 设球半径为 \(R\)，球心到截面的距离为 \(R/2\)，则由截面半径关系 \(\left(\frac{2}{\sqrt{3}}\right)^2=R^2-\left(\frac{R}{2}\right)^2=\frac{3R^2}{4}\). 解得 \(R^2=16/9\)，所以球面面积为 \(4\pi R^2=\frac{64}{9}\pi\).
故选 D.
\end{solution}

\begin{problem}
函数 \(y = \arccos (\sin x)\)（\(-\frac{\pi}{3} < x < \frac{2\pi}{3}\)）的值域是（\quad）

\choices
    {\(\left(\frac{\pi}{6}, \frac{5\pi}{6}\right)\)}
    {\(\left[0, \frac{5\pi}{6}\right)\)}
    {\(\left(\frac{\pi}{3}, \frac{2\pi}{3}\right)\)}
    {\(\left[\frac{\pi}{6}, \frac{2\pi}{3}\right)\)}
\end{problem}

\begin{answer}
B
\end{answer}

\begin{solution}
在区间 \(\left(-\pi/3,2\pi/3\right)\) 内，\(\sin x\) 的最大值 \(1\) 在 \(x=\pi/2\) 处取得，最小值趋近于 \(-\sqrt{3}/2\) 但不能取得，所以 \(\sin x\in\left(-\frac{\sqrt{3}}{2},1\right]\).
因为 \(\arccos t\) 在 \([-1,1]\) 上递减，故
\(\arccos(\sin x)\in\left[0,\arccos\left(-\frac{\sqrt{3}}{2}\right)\right)=\left[0,\frac{5\pi}{6}\right)\).
故选 B.
\end{solution}

\begin{problem}
定义在 \((- \infty, + \infty)\) 上的任意函数 \(f(x)\) 都可以表示成一个奇函数 \(g(x)\) 和一个偶函数 \(h(x)\) 之和，如果 \(f(x) = \lg \left(10^x + 1\right)\)，\(x \in (-\infty, +\infty)\)，那么（\quad）

\choices
    {\(g(x) = x\)，\(h(x) = \lg \left(10^x + 10^{-x} + 2\right)\)}
    {\(g(x) = \frac{1}{2}\left[\lg \left(10^x + 1\right) + x\right]\)，\(h(x) = \frac{1}{2}\left[\lg \left(10^x + 1\right) - x\right]\)}
    {\(g(x) = \frac{x}{2}\)，\(h(x) = \lg \left(10^x + 1\right) - \frac{x}{2}\)}
    {\(g(x) = -\frac{x}{2}\)，\(h(x) = \lg \left(10^x + 1\right) + \frac{x}{2}\)}
\end{problem}

\begin{answer}
C
\end{answer}

\begin{solution}
记 \(f(x)=\lg(10^x+1)\). 由
\[
f(-x)=\lg(10^{-x}+1)=\lg\frac{10^x+1}{10^x}=f(x)-x
\]
可得奇函数部分和偶函数部分分别为 \(g(x)=\frac{f(x)-f(-x)}{2}=\frac{x}{2}\)，\(h(x)=\frac{f(x)+f(-x)}{2}=f(x)-\frac{x}{2}\).
所以选项 C 正确.
\end{solution}

\section{填空题}

\begin{problem}
在 \((3 - x)^7\) 的展开式中，\(x^5\) 的系数是\fillinblank{}.
\end{problem}

\begin{answer}
\(-189\)
\end{answer}

\begin{solution}
二项式展开式的通项为 \(\mathrm C_{7}^{k}3^{7-k}(-x)^k\). 取 \(k=5\)，得到 \(x^5\) 的系数 \(\mathrm C_{7}^{5}3^2(-1)^5=21\cdot9\cdot(-1)=-189\).
\end{solution}

\begin{problem}
抛物线 \(y^{2} = 8 - 4x \) 的准线方程是\fillinblank{}，圆心在该抛物线的顶点且与其准线相切的圆的方程是\fillinblank{}.
\end{problem}

\begin{answer}
\(x=3\)，\(\left(x-2\right)^2+y^2=1\)
\end{answer}

\begin{solution}
将抛物线方程写成 \(y^2=-4(x-2)\). 与标准形式 \(y^2=4p(x-2)\) 比较，得 \(p=-1\)，所以准线为 \(x=2-p=3\). 抛物线顶点为 \((2,0)\). 以该点为圆心且与直线 \(x=3\) 相切的圆半径为 \(1\)，故圆的方程为 \((x-2)^2+y^2=1\).
\end{solution}

\begin{problem}
已知 \(\sin \theta + \cos \theta = \frac{1}{5}\)，\(\theta \in (0, \pi)\)，则 \(\cot \theta\) 的值是\fillinblank{}.
\end{problem}

\begin{answer}
\(-\frac{3}{4}\)
\end{answer}

\begin{solution}
由 \(\sin\theta+\cos\theta=1/5\)，得 \(1+2\sin\theta\cos\theta=\frac{1}{25}\)，即 \(\sin\theta\cos\theta=-12/25\). 因 \(\theta\in(0,\pi)\)，有 \(\sin\theta>0\)；结合 \(\sin\theta\cos\theta<0\)，得 \(\cos\theta<0\)，从而 \(\sin\theta-\cos\theta>0\). 于是 \((\sin\theta-\cos\theta)^2=1-2\sin\theta\cos\theta=\frac{49}{25}\)，所以 \(\sin\theta-\cos\theta=7/5\)，并且 \(\sin\theta+\cos\theta=\frac15\).
两式相加、相减，得 \(\sin\theta=4/5\)，\(\cos\theta=-3/5\). 因此
\(\cot\theta=\frac{\cos\theta}{\sin\theta}=-\frac34\).
\end{solution}

\begin{problem}
设圆锥底面圆周上两点 \(A\)，\(B\) 间的距离为 \(2\)，圆锥顶点到直线 \(AB\) 的距离为 \(\sqrt{3} \)，\(AB\) 和圆锥的轴的距离为 \(1\)，则该圆锥的体积为\fillinblank{}.
\end{problem}

\begin{answer}
\(\frac{2\sqrt{2}}{3}\pi\)
\end{answer}

\begin{solution}
设圆锥底面圆心为 \(O\)，\(M\) 为弦 \(AB\) 的中点，顶点为 \(S\). 由圆心到弦中点的连线垂直于弦，得 \(OM\perp AB\)，且 \(AM=1\)，\(OM=1\).
在直角三角形 \(OAM\) 中，底面半径满足
\[
r^2=OA^2=OM^2+AM^2=2
\]
又 \(SO\) 是圆锥的轴，\(SO\perp\) 底面，故 \(SO\perp AB\). 因为 \(\overrightarrow{SM}=\overrightarrow{SO}+\overrightarrow{OM}\)，且 \(SO\perp AB\)、\(OM\perp AB\)，所以 \(SM\perp AB\)，故 \(SM\) 就是题设给出的顶点到直线 \(AB\) 的距离，且 \(SM^2=SO^2+OM^2\).
由题意 \(SM=\sqrt{3}\)，得 \(SO^2=3-1=2\)，所以圆锥高 \(h=\sqrt{2}\). 因此圆锥体积为
\[
V=\frac13\pi r^2h=\frac13\pi\cdot2\cdot\sqrt{2}
=\frac{2\sqrt{2}}{3}\pi
\]
\end{solution}

\begin{problem}
在测量某物理量的过程中，因仪器和观察的误差，使得 \(n\) 次测量分别得到 \(a_{1}\)，\(a_{2}\)，\(\cdots\)，\(a_{n}\)，共 \(n\) 个数据，我们规定所测量物理量的“最佳近似值”\(a\) 是这样一个量：与其他近似值比较，\(a\) 与各数据的差的平方和最小. 依此规定，从 \(a_{1}\)，\(a_{2}\)，\(\cdots\)，\(a_{n}\) 推出的 \(a =\fillinblank{}\).
\end{problem}

\begin{answer}
\(\frac{a_1+a_2+\cdots+a_n}{n}\)
\end{answer}

\begin{solution}
设近似值为 \(t\)，考察误差平方和 \(F(t)=\sum_{i=1}^{n}(t-a_i)^2\). 令 \(\bar a=\frac{a_1+a_2+\cdots+a_n}{n}\).
由于 \(\sum_{i=1}^{n}(\bar a-a_i)=0\)，有
\[
F(t)-F(\bar a)
=\sum_{i=1}^{n}\left((t-a_i)^2-(\bar a-a_i)^2\right)
=n(t-\bar a)^2\geqslant0
\]
所以 \(F(t)\) 在 \(t=\bar a\) 时取得最小值，最佳近似值为 \(a=\frac{a_1+a_2+\cdots+a_n}{n}\).
\end{solution}

\section{解答题}

\begin{problem}
已知 \(z = 1 + \i\).

\begin{enumerate}
\item 设 \(\omega = z^2 + 3\overline{z} - 4\)，求 \(\omega\) 的三角形式；
\item 如果 \(\frac{z^{2}+az+b}{z^{2}-z+1}=1-\i\)，求实数 \(a\)，\(b\) 的值.
\end{enumerate}
\end{problem}

\begin{solution}
\begin{enumerate}
\item 由 \(z=1+\i\) 得 \(\overline{z}=1-\i\)，\(z^2=2\i\). 因此
\[
\omega=z^2+3\overline{z}-4
=2\i+3(1-\i)-4=-1-\i
\]
于是 \(|\omega|=\sqrt{2}\)，且 \(\omega\) 在第三象限，所以其辐角可取 \(5\pi/4\). 故
\[
\omega=\sqrt{2}\left(\cos\frac{5\pi}{4}+\i\sin\frac{5\pi}{4}\right)
\]
\item 先计算分母
\(z^2-z+1=2\i-(1+\i)+1=\i\neq0\).
分子为 \(z^2+az+b=2\i+a(1+\i)+b=(a+b)+(a+2)\i\). 由题设，分子等于 \(\i(1-\i)=1+\i\). 根据复数相等的定义，得 \(a+b=1\)，\(a+2=1\).
解得 \(a=-1\)，\(b=2\).
\end{enumerate}
\end{solution}

\begin{problem}
已知函数 \(f(x) = \tan x\)，\(x \in \left(0, \frac{\pi}{2}\right)\). 若 \(x_1\)，\(x_2 \in \left(0, \frac{\pi}{2}\right)\)，且 \(x_1 \neq x_2 \)，证明：\(\frac{1}{2}\left[f(x_1) + f(x_2)\right] > f\left(\frac{x_1 + x_2}{2}\right) \).
\end{problem}

\begin{solution}
令 \(m=\frac{x_1+x_2}{2}\)，\(d=\frac{x_1-x_2}{2}\).
则 \(m\in(0,\pi/2)\)，且 \(d\neq0\). 由 \(x_1\)，\(x_2\in(0,\pi/2)\) 可知 \(\cos x_1>0\)、\(\cos x_2>0\)，并且
\[
\tan x_1+\tan x_2
=\frac{\sin(x_1+x_2)}{\cos x_1\cos x_2}
=\frac{\sin 2m}{\cos x_1\cos x_2}
\]
又
\[
\cos x_1\cos x_2
=\cos(m+d)\cos(m-d)
=\cos^2m-\sin^2d<\cos^2m
\]
因为 \(0<m<\pi/2\)，所以 \(\sin2m>0\). 因此
\[
\frac12\left(\tan x_1+\tan x_2\right)-\tan m
=\frac{\sin2m}{2}\left(\frac{1}{\cos x_1\cos x_2}-\frac{1}{\cos^2m}\right)
=\frac{\sin2m\sin^2d}{2\cos x_1\cos x_2\cos^2m}>0
\]
于是 \(\frac12\left[f(x_1)+f(x_2)\right]>f\left(\frac{x_1+x_2}{2}\right)\).
\end{solution}

\begin{problem}
如图，已知 \(A_{1}B_{1}C_{1} - ABC\) 是正三棱柱，\(D\) 是 \(AC\) 中点.

\begin{enumerate}
\item 证明：\(AB_{1} \parallel \text{平面 } DBC_{1}\)；
\item 假设 \(AB_{1} \perp BC_{1}\)，求以 \(BC_{1}\) 为棱，\(DBC_{1}\) 与 \(CBC_{1}\) 为面的二面角 \(\alpha\) 的度数.
\end{enumerate}

\bitmapfigure[trim=0 5.29cm 0 0,clip,max width=0.42\linewidth]{img/1994/national_science/q23_fig1.png}
\end{problem}

\begin{solution}
\begin{enumerate}
\item 在矩形 \(B_1BCC_1\) 中，对角线 \(B_1C\) 与 \(BC_1\) 互相平分. 设它们的交点为 \(E\)，则 \(E\) 是 \(B_1C\) 的中点. 又 \(D\) 是 \(AC\) 的中点，所以在三角形 \(AB_1C\) 中，中位线定理给出 \(DE\parallel AB_1\).
由于 \(D\)，\(E\in\text{平面 }DBC_1\)，有 \(DE\subset\text{平面 }DBC_1\). 又底面 \(ABC\) 与平面 \(DBC_1\) 的交线为 \(BD\)，而 \(A\notin BD\)，所以 \(A\notin\text{平面 }DBC_1\). 从而
\(AB_1\parallel\text{平面 }DBC_1\).
\item 设正三角形 \(ABC\) 的边长为 \(a\)，棱柱高为 \(h\). 取空间直角坐标系，使
\(B=(0,0,0)\)，\(C=(a,0,0)\)，\(A=\left(\frac a2,\frac{\sqrt3a}{2},0\right)\)，\(B_1=(0,0,h)\)，\(C_1=(a,0,h)\). 则 \(\overrightarrow{AB_1}=\left(-\frac a2,-\frac{\sqrt3a}{2},h\right)\)，\(\overrightarrow{BC_1}=(a,0,h)\).
由 \(AB_1\perp BC_1\)，得 \(\overrightarrow{AB_1}\cdot\overrightarrow{BC_1}=-\frac{a^2}{2}+h^2=0\)，从而 \(h^2=\frac{a^2}{2}\).
点 \(D\) 的坐标为 \(D=\left(\frac{3a}{4},\frac{\sqrt3a}{4},0\right)\).
平面 \(DBC_1\) 内的两个相交方向向量可取 \(\overrightarrow{BC_1}\) 和 \(\overrightarrow{BD}\)，故其法向量可取
\(\overrightarrow{BC_1}\times\overrightarrow{BD}=\frac a4\left(-\sqrt3h,3h,\sqrt3a\right)\).
平面 \(CBC_1\) 为 \(y=0\)，可取法向量 \(\bs{n}_2=(0,1,0)\). 两平面所成的锐二面角为 \(\alpha\)，于是
\[
\cos\alpha
=\frac{3h}{\sqrt{12h^2+3a^2}}
=\frac{3a/\sqrt2}{3a}
=\frac{1}{\sqrt2}
\]
因此 \(\alpha=45^\circ\).
\end{enumerate}
\end{solution}

\begin{problem}
已知直线 \(l\) 过坐标原点，抛物线 \(C\) 顶点在原点，焦点在 \(x\) 轴正半轴上. 若点 \(A(-1,0)\) 和点 \(B(0,8)\) 关于 \(l\) 的对称点都在 \(C\) 上，求直线 \(l\) 和抛物线 \(C\) 的方程.
\end{problem}

\begin{solution}
设抛物线 \(C\) 的方程为 \(y^2=2px\)，其中 \(p>0\). 直线 \(l\) 不能是坐标轴：若 \(l\) 为 \(x\) 轴，则点 \(A\) 的对称点仍为 \(A\)，其横坐标为负，不在 \(C\) 上；若 \(l\) 为 \(y\) 轴，则点 \(B\) 的对称点仍为 \(B\)，而 \(C\) 在 \(x=0\) 时只有原点. 因此可设
\(l:y=kx\)，\(k\neq0\).
设 \(X=(x,y)\)，直线 \(l\) 的方向向量为 \((1,k)\). 点 \(X\) 在 \(l\) 上的正投影为
\(\frac{x+ky}{1+k^2}(1,k)\).
反射点的位矢等于投影点位矢的两倍减去 \(X\) 的位矢，所以关于直线 \(y=kx\) 的对称变换为
\[
(x,y)\longmapsto
\left(\frac{(1-k^2)x+2ky}{1+k^2},
\frac{2kx+(k^2-1)y}{1+k^2}\right)
\]
所以 \(A'=\left(\frac{k^2-1}{1+k^2},-\frac{2k}{1+k^2}\right)\)，\(B'=\left(\frac{16k}{1+k^2},\frac{8(k^2-1)}{1+k^2}\right)\).
由于 \(A'\)、\(B'\) 在开口向右的抛物线上，二者横坐标均为正，故 \(k>0\) 且 \(k>1\). 将 \(A'\) 代入抛物线方程，得
\(p=\frac{2k^2}{(1+k^2)(k^2-1)}\).
将 \(B'\) 代入，得
\(p=\frac{2(k^2-1)^2}{k(1+k^2)}\).
两式相等，且 \(k>0\)，从而
\(\frac{k^2}{k^2-1}=\frac{(k^2-1)^2}{k}\)，\(k^3=(k^2-1)^3\)
所以 \(k=k^2-1\). 结合 \(k>1\)，得 \(k=\frac{1+\sqrt5}{2}\).
又 \(k^2-1=k\)，代入 \(p\) 的表达式，得
\(p=\frac{2k}{1+k^2}=\frac{2}{\sqrt5}\). 因此 \(l:y=\frac{1+\sqrt5}{2}x\)，\(C:y^2=2px=\frac{4}{\sqrt5}x=\frac{4\sqrt5}{5}x\).
\end{solution}

\begin{problem}
设 \(\{a_{n}\}\) 是正数组成的数列，其前 \(n\) 项和为 \(S_{n}\)，且对于所有的自然数 \(n\)，\(a_{n}\) 与 \(2\) 的等差中项等于 \(S_{n}\) 与 \(2\) 的等比中项.

\begin{enumerate}
\item 写出数列 \(\{a_{n}\}\) 的前 \(3\) 项；
\item 求数列 \(\{a_{n}\}\) 的通项公式（写出推证过程）；
    \item 令 \(b_{n} = \frac{1}{2}\left(\frac{a_{n+1}}{a_{n}} + \frac{a_{n}}{a_{n+1}}\right)\)（\(n \in \mathbb{N}\)），求 \(\displaystyle\lim_{n \to \infty}\left(b_{1} + b_{2} + \cdots + b_{n} - n\right)\).
\end{enumerate}
\end{problem}

\begin{solution}
\begin{enumerate}
\item “\(a_n\) 与 \(2\) 的等差中项”等于 \((a_n+2)/2\)，“\(S_n\) 与 \(2\) 的等比中项”等于 \(\sqrt{2S_n}\). 由题意 \(\frac{a_n+2}{2}=\sqrt{2S_n}\)，且两边均为正数，平方得 \(S_n=\frac{(a_n+2)^2}{8}\).
当 \(n=1\) 时，\(S_1=a_1\)，所以
\(8a_1=(a_1+2)^2\)，即 \((a_1-2)^2=0\)，即 \(a_1=2\). 当 \(n=2\) 时，\(S_2=a_1+a_2=2+a_2\)，故 \(8(2+a_2)=(a_2+2)^2\)，
即 \(a_2=6\)，因为另一根为负数. 当 \(n=3\) 时，\(S_3=8+a_3\)，故
\(8(8+a_3)=(a_3+2)^2\)，
即 \((a_3-10)(a_3+6)=0\). 由 \(a_3>0\)，得 \(a_3=10\).
\item 对任意 \(n\)，由上面的关系分别写出 \(S_n\) 与 \(S_{n+1}\)，并利用 \(S_{n+1}-S_n=a_{n+1}\)，得
\[
(a_{n+1}+2)^2-(a_n+2)^2=8a_{n+1}
\]
整理为
\((a_{n+1}+a_n)(a_{n+1}-a_n-4)=0\).
由于数列各项为正数，\(a_{n+1}+a_n>0\)，所以
\(a_{n+1}-a_n=4\).
故 \(\{a_n\}\) 是首项 \(2\)、公差 \(4\) 的等差数列，从而
\(a_n=2+4(n-1)=4n-2\).
\item 由通项公式，
\[
b_n=\frac12\left(\frac{4n+2}{4n-2}+\frac{4n-2}{4n+2}\right)
=\frac{4n^2+1}{4n^2-1}
=1+\frac{2}{4n^2-1}
\]
因此
\[
b_1+\cdots+b_n-n
=\sum_{k=1}^{n}\frac{2}{4k^2-1}
=\sum_{k=1}^{n}\left(\frac{1}{2k-1}-\frac{1}{2k+1}\right)
=1-\frac{1}{2n+1}
\]
令 \(n\to\infty\)，得到 \(\lim_{n\to\infty}\left(b_1+b_2+\cdots+b_n-n\right)=1\).
\end{enumerate}
\end{solution}
